%! Equivalent Representations of Polynomial and Rational Expressions — Unit 1, Lesson 1.11 — Exit Ticket (Answer Key)
\documentclass[10pt]{article}
\usepackage{precalculus-article}
\usepackage{precalculus-key}   % pulls in -boxes; do NOT also load -boxes
\newcommand{\TallMath}[1]{$\displaystyle #1\rule[-1.4em]{0pt}{3.2em}$}

\begin{document}

\pageheader{Unit 1, Lesson 1.11}{Exit Ticket --- Answer Key}
\namedateperiod

\vspace{0.14in}

\begin{enumerate}[itemsep=16pt]

\item A student needs to find the \textbf{vertical asymptotes and holes} of a rational function.
      Circle the form that is most useful, and explain why in one phrase.

      \medskip
      \ans{A. Factored form} \qquad\qquad \textbf{B.} Standard form

      \vspace{0.06in}
      Because: \ansline{factored form shows the zeros of the numerator and denominator, making it easy to identify holes and vertical asymptotes by comparing shared factors.}

\item Perform the first step of polynomial long division to divide $x^2 + 3x + 5$ by $x + 2$.

      What is the first term of the quotient? \ans{$x$}

      After multiplying and subtracting, what is the new dividend? \ans{$x + 5$}

\item Use Pascal's Triangle to expand $(x + 2)^4$.
      (Row 4 of Pascal's Triangle: $1 \quad 4 \quad 6 \quad 4 \quad 1$)

      \vspace{0.06in}
      $(x+2)^4 = $ \ans{$x^4 + 8x^3 + 24x^2 + 32x + 16$}

\end{enumerate}

\begin{teachernote}
Item 2: First term of quotient $= x^2 \div x = x$. Multiply: $x(x+2) = x^2+2x$. Subtract from $x^2+3x+5$: new dividend is $x+5$. Final answer (if completed): quotient $x+1$, remainder $3$, so $\frac{x^2+3x+5}{x+2} = x+1+\frac{3}{x+2}$.

Item 3: $(x+2)^4 = 1\cdot x^4 + 4\cdot x^3\cdot 2 + 6\cdot x^2\cdot 4 + 4\cdot x\cdot 8 + 1\cdot 16 = x^4+8x^3+24x^2+32x+16$.
\end{teachernote}

\end{document}
